%Document-Author: Bonato Paolo + Bortolazzo Matteo + Maino Elia
%Document-Date: 2016/01/13
%Document-Description: Documento di Analisi dei requisiti del gruppo SWEeneyThreads 

\documentclass[a4paper]{article}
\usepackage[english, italian]{babel}
\usepackage[T1]{fontenc}
\usepackage[utf8]{inputenc}
\usepackage{url}
\usepackage{graphicx}
\usepackage[hidelinks]{hyperref}
\usepackage{booktabs}
\usepackage{eurosym}
\usepackage{tabularx}
\usepackage{pifont}
\usepackage[table]{xcolor}
\usepackage{float}
\usepackage[]{appendix}
\usepackage{ltxtable} 
\usepackage{geometry}
\geometry{margin=1in}
\usepackage{longtable}
\usepackage{multirow}

\graphicspath{{Immagini/}}

\newcolumntype{Y}{>{\centering\arraybackslash}X}
\newcolumntype{s}{>{\hsize=.21\hsize}X}
\newcolumntype{f}{>{\hsize=.37\hsize}X}
\newcolumntype{m}{>{\hsize=.42\hsize}X}
\newcolumntype{t}{>{\hsize=.1\hsize}X}
\newcolumntype{r}{>{\hsize=.3\hsize}X}
\newcolumntype{k}{>{\hsize=.4\hsize}X}

\renewcommand{\abstractname}{Tabella contenuti}

\begin{document}
	
	\begin{titlepage}
		% Defines a new command for the horizontal lines, change thickness here
		\newcommand{\HRule}{\rule{\linewidth}{0.5mm}} 
		\center  
		
		% HEADING SECTION
		\textsc{\LARGE SweeneyThreads}\\[1.5cm] 
		\textsc{\Large Actorbase}\\[0.5cm] 
		\textsc{\large a NoSQL DB based on the Actor model}\\[0.5cm]
		
		
		% TITLE SECTION
		\HRule \\[0.4cm]
		{ \huge \bfseries Analisi dei requisiti}\\[0.4cm] 
		\HRule \\[1.5cm]
		
		% AUTHOR SECTION
		\begin{minipage}{0.4\textwidth}
			\begin{flushleft} \large
				\emph{Redattori:}\\
				Bonato Paolo \\
				Bortolazzo Matteo \\
				Maino Elia
			\end{flushleft}
		\end{minipage}
		~
		\begin{minipage}{0.4\textwidth}
			\begin{flushright} \large
				\emph{Approvazione:} \\
				Padovan Tommaso \\
				\emph{Verifica:} \\
				Biggeri Mattia \\
				Tommasin Davide 
			\end{flushright}
		\end{minipage}
		
		%immagine
		\begin{figure}[H]
			\centering
			\includegraphics[scale=0.8]{sweeney.png}
		\end{figure}
		\begin{center}
			Versione 1.2.1
		\end{center}
		% Date, change the \today to a set date if you want to be precise
		{\large \today}\\[3cm] 
		% Fill the rest of the page with whitespace
		\vfill  
	\end{titlepage}
	
	
	\tableofcontents
	
	\newpage \section*{Diario delle modifiche}
		\begin{table}[H]
			\begin{tabularx}{\textwidth}{s f m X}
				\noalign{\hrule height 1.5pt}
				\rowcolor{orange!85} Versione & Data & Autore & Descrizione \\
				\noalign{\hrule height 0.5pt}
				1.2.1 & 2016-03-05 & \emph{Analista} \newline Maino Elia & Ridefinizione dei casi d'uso: stesura di diagrammi e descrizioni fino a UC 3.1
				 \\
				\noalign{\hrule height 0.5pt}
				1.2.0 & 2016-01-21 & \emph{Responsabile} \newline Padovan Tommaso  & Approvazione documento
				 \\
				 \noalign{\hrule height 0.5pt}
				1.1.1 & 2016-01-21 & \emph{Analisti} \newline Bonato Paolo \newline Bortolazzo Matteo\newline Maino Elia  & 
				Correzioni degli errori individuati nella verifica
				 \\
				 \noalign{\hrule height 0.5pt}
				1.1.0 & 2016-01-21 & \emph{Verificatori} \newline Biggeri Mattia \newline Tommasin Davide  & Verifica del documento \\
				 \noalign{\hrule height 0.5pt}
				1.0.5 & 2016-01-20 & \emph{Analista} \newline Bortolazzo Matteo  & Stesura lista dei requisiti dedotti dal capitolato e da analisi interne
				 \\
				 \noalign{\hrule height 0.5pt}
				1.0.4 & 2016-01-19 & \emph{Analisti} \newline Bonato Paolo \newline Maino Elia & Stesura lista dei requisiti dedotti dai casi d'uso
				 \\
				\noalign{\hrule height 0.5pt}
				1.0.3 & 2016-01-18 & \emph{Analista} \newline Maino Elia & Stesura dei casi d'uso (da 4.x a 5.x) \\
				\noalign{\hrule height 0.5pt}
				1.0.2 & 2016-01-17 & \emph{Analista} \newline Maino Elia & Stesura dei casi d'uso (da 1.x a 3.x) \\
				\noalign{\hrule height 0.5pt}
				1.0.1 & 2016-01-17 & \emph{Analista} \newline Maino Elia & Stesura caratteristiche generali del prodotto \\
				\noalign{\hrule height 0.5pt}
				1.0.0 & 2016-01-17 & \emph{Analista} \newline Maino Elia & Creazione scheletro documento e stesura introduzione \\
				\noalign{\hrule height 1.5pt}
			\end{tabularx}
			\caption{Diario delle modifiche \label{tab:table_label}}
		\end{table}
	



	\newpage \section{Introduzione}
	\subsection{Scopo del documento}
		Lo scopo del seguente documento è presentare l'insieme di requisiti individuati dal gruppo 
		SWEeneyThreads durante l'analisi del \emph{Capitolato C1}, gli incontri con il committente 
		\emph{Cardin Riccardo} e l'analisi dei casi d'uso. 
		
		Si intende fornire una rappresentazione ordinata delle funzionalità che il prodotto \emph{Actorbase} 
		offrirà al momento del rilascio.
	\subsection{Scopo del prodotto}
		Lo scopo del progetto è la realizzazione di un DataBase NoSQL key-value basato sul modello ad 
		Attori con l'obiettivo di fornire una tecnologia adatta allo sviluppo di moderne 
		applicazioni che richiedono brevissimi tempi di risposta e che elaborano enormi quantità 
		di dati. Lo sviluppo porterà al rilascio del software sotto licenza MIT.
	\subsection{Glossario}
		Con lo scopo di evitare ambiguità di linguaggio e di massimizzare la comprensione dei documenti, il 
      gruppo ha steso un documento interno che è il \emph{Glossario v1.0.3}. In esso saranno definiti, in modo
      chiaro e conciso i termini che possono causare ambiguità o incomprensione del testo.
	\subsection{Riferimenti}
	\subsubsection{Informativi}	
		\begin{itemize}
			\item \textbf{Slide dell'insegnamento Ingegneria del software mod.A:} \\
			\url{http://www.math.unipd.it/~tullio/IS-1/2015/Dispense/E02.pdf}
		\end{itemize}
	\subsubsection{Normativi}
		\begin{itemize}
			\item \textbf{Norme di progetto:} \emph{Norme di progetto v1.1.2}
			\item \textbf{Capitolato d'appalto Actorbase (C1):} \\ 
			\url{http://www.math.unipd.it/~tullio/IS-1/2015/Progetto/C1p.pdf}
		\end{itemize}
		
		
	\newpage \section{Caratteristiche generali del prodotto}
	\subsection{Obiettivi del prodotto}
		Il prodotto si pone l'obiettivo di fornire un database NoSQL basato sul modello ad attori che possa 
		essere utilizzato con successo dalle cosiddette \emph{Reactive Applications}, ovvero applicazioni
		orientate agli eventi, scalabili, resilienti e responsive. Per tali applicazioni il paradigma di accesso 
		ai dati basato su database di tipo relazionale risulta non applicabile, poichè troppo limitante, sono 
		dunque necessarie nuove forme di gestione dell'informazione, \emph{Actorbase} intende proporsi 
		come una valida opzione.
	\subsection{Funzioni principali del prodotto}
		\emph{Actorbase} prevede un'interazione con l'utente da riga di comando attraverso l'uso di un 
		\emph{Domain Specific Language (DSL)}. 
		\\ \\
		Una volta effettuata la connessione al server, l'utente potrà:
		\begin{itemize}
			\item Visualizzare i database presenti sul server
			\item Creare un nuovo database
			\item Selezionare un database da utilizzare
			\item Rinominare un database già presente
			\item Eliminare un database
			\item Esportare i database (backup)
		\end{itemize}
		La selezione di un database permetterà all'utente di effettuare una serie di operazioni di interrogazione e modifica di esso:
		\begin{itemize}
			\item Creazione di mappe
			\item Selezione di una mappa
			\item Rimozione di una mappa
			\item Visualizzazione delle mappe presenti
		\end{itemize}
		La selezione di una mappa permetterà all'utente di accedere alle operazioni di interrogazione e modifica a tale livello:
		\begin{itemize}
			\item Inserimento di un item
			\item Ricerca di un item per chiave
			\item Cancellazione di un item
			\item Aggiornamento di un item già presente
		\end{itemize}
	\subsection{Target d'utenza}
		Il prodotto si rivolge a sviluppatori di applicazioni moderne, che trattano enormi moli di dati (nell'
		ordine dei Petabyte), che richiedono brevissimi tempi di risposta e che necessitano di un uptime 
		del 100\%.
	\subsection{Vincoli per l'utilizzo}
		L'esecuzione di \emph{Actorbase} avviene attraverso la JVM (Java Virtual Machine), il corretto funzionamento di tutte le caratteristiche del prodotto è garantito su macchine che utilizzano la versione 8 o successive della JVM. 
	
	\newpage \section{Casi d'uso}
		Di seguito viene riportata la descrizione accurata di tutti i casi d'uso individuati dal gruppo a seguito
		delle seguenti attività:
		\begin{itemize}
			\item Analisi del capitolato C1 Actorbase
			\item Confronti con il committente \emph{Cardin Riccardo}
			\item Riunioni e discussioni interne al gruppo
			\item Analisi della struttura e delle funzionalità di altri database non realzionali
		\end{itemize}
		La struttura di un caso d'uso è definita nel documento \emph{Norme di progetto v1.1.1 sez 2.1.3}.
		\subsection{Attori}
		L'interazione degli utenti con l'applicativo avverrà tramite una CLI (Command Line Interface).
		Sono state individuate due tipologie principali di attori che interagiscono con il sistema:
		\begin{itemize}
			\item \textbf{Utente non autenticato:} è l'utente che ha avviato l'applicativo e interagisce con la CLI senza aver ancora effettuato con successo l'operazione di connessione al server;
			\item \textbf{Utente autenticato:} è l'utente che ha già effettuato con successo l'operazione di connessione ad un server che utilizza \emph{Actorbase}. 
		\end{itemize}
		\emph{Actorbase} non prevede un meccanismo di permessi per quanto riguarda i dati presenti su un server. Una volta che l'utente si è autenticato con successo ha la possibilità di effettuare tutte le operazioni disponibili. Per tale motivo non è presente un utente amministratore, in quanto l'amministrazione del server è già una funzionalità garantita all'utente autenticato.
		\subsection{Visione ad alto livello delle interazioni tramite CLI}
		 	\begin{figure}[H]
				\centering
				\includegraphics[scale=0.2]{UC/"Actorbase".png}
				\caption{Interazioni tramite CLI}
			\end{figure}
			Il diagramma in figura illustra le principali operazioni che un utente (autenticato e non) può
			 effettuare sul sistema tramite l'interfaccia da riga di comando. 
			\begin{itemize}
				\item Connessione al server
				\item Visualizzazione di un aiuto
				\item Operazioni a livello di database
				\item Operazioni a livello di mappa
				\item Operazioni a livello di riga
				\item Disconnessione dal server
			\end{itemize}
	In ogni momento un utente può inserire un comando inesistente o non utilizzabile in quel determinato momento. In tal caso nessuna azione viene eseguita dal programma e l'utente riceve un messaggio di errore esplicativo seguito dalla possibilità di inserire un nuovo comando.	 
	 
	 
	 
	 \subsection{UC 1: Connessione al server}
	 \begin{figure}[H]
				\centering
				\includegraphics[scale=0.3]{UC/"UC 1 Connessione".png}
				\caption{Diagramma di UC1: Connessione al server}
			\end{figure}
	\textbf{Descrizione} 
	\\ \\
	L'utente ha appena avviato l'interfaccia CLI dell'applicativo e intende connettersi ad un server contenente i database utilizzando il comando \texttt{CONNECT} (UC 1.1). L'utente deve inserire l'indirizzo del server (UC 1.2) e successivamente l'username e la password necessari per l'accesso (UC 1.3), nel caso in cui il tentativo di connessione dovesse fallire, l'utente riceve un messaggio di errore informativo.
	\begin{table}[H]
			\begin{tabularx}{\textwidth}{r X}
				\textbf{Codice gerarchico} & UC1 \\
				\noalign{\hrule height 0.5pt}
				\textbf{Nome sintetico} & Connessione al server \\
				\noalign{\hrule height 0.5pt}
				\textbf{Attore principale} & Utente non autenticato\\
				\noalign{\hrule height 0.5pt}
				\textbf{Attori secondari} & Nessuno \\
				\noalign{\hrule height 0.5pt}
				\textbf{Pre-condizione} & L'utente ha avviato l'interfaccia CLI, non è autenticato e intende connettersi ad un server\\
				\noalign{\hrule height 0.5pt}
				\textbf{Post-condizione} & L'utente risulta connesso al server \\
				\noalign{\hrule height 0.5pt}
				\textbf{Flusso eventi} & \begin{enumerate}
				\item L'utente inserisce il comando \texttt{CONNECT} (UC 1.1)
				\item L'utente inserisce l'indirizzo del server (UC 1.2)
				\item L'utente inserisce username e password (UC 1.3)
				\end{enumerate} \\
				\noalign{\hrule height 0.5pt}
				\textbf{Scenari alternativi} & Nessuno \\
				\noalign{\hrule height 0.5pt}
				\textbf{Lista requisiti\newline dedotti} & \begin{itemize}
				\item ...
				\end{itemize} 
			\end{tabularx}
			\caption{Caso d'uso UC 1 - Connessione al server}
		 \end{table} 
	 
	 
	\subsection{UC 1.1: Inserimento comando \texttt{CONNECT} }
	 \textbf{Descrizione}
	 \\ \\
	 L'utente intende effettuare una connessione al server, deve avviare la procedura di connessione inserendo il comando \texttt{CONNECT} da terminale.
	\begin{table}[H]
			\begin{tabularx}{\textwidth}{r  X}
				\textbf{Codice gerarchico} & UC1.1 \\
				\noalign{\hrule height 0.5pt}
				\textbf{Nome sintetico} & Inserimento comando \texttt{CONNECT} \\
				\noalign{\hrule height 0.5pt}
				\textbf{Attore principale} & Utente non autenticato\\
				\noalign{\hrule height 0.5pt}
				\textbf{Attori secondari} & Nessuno \\
				\noalign{\hrule height 0.5pt}
				\textbf{Pre-condizione} & L'utente intende avviare la procedura di connessione al server\\
				\noalign{\hrule height 0.5pt}
				\textbf{Post-condizione} & L'utente ha inserito con successo il comando di connessione\\
				\noalign{\hrule height 0.5pt}
				\textbf{Flusso eventi} & \begin{enumerate}
				\item L'utente scrive sul terminale il comando di connessione
				\end{enumerate} \\
				\noalign{\hrule height 0.5pt}
				\textbf{Scenari alternativi} & Nessuno \\
				\noalign{\hrule height 0.5pt}
				\textbf{Lista requisiti\newline dedotti} & \begin{itemize}
				\item ...
				\end{itemize} 
			\end{tabularx}
			\caption{Caso d'uso UC 1.1 - Inserimento comando \texttt{CONNECT}}
		 \end{table} 	 
	 
	 
	 \subsection{UC 1.2: Inserimento indirizzo server }
	 \textbf{Descrizione}
	 \\ \\
	 L'utente sta effettuando l'operazione di connessione al server e deve inserire l'indirizzo del server a cui vuole connettersi.
	\begin{table}[H]
			\begin{tabularx}{\textwidth}{r  X}
				\textbf{Codice gerarchico} & UC1.2 \\
				\noalign{\hrule height 0.5pt}
				\textbf{Nome sintetico} & Inserimento indirizzo server \\
				\noalign{\hrule height 0.5pt}
				\textbf{Attore principale} & Utente non autenticato\\
				\noalign{\hrule height 0.5pt}
				\textbf{Attori secondari} & Nessuno \\
				\noalign{\hrule height 0.5pt}
				\textbf{Pre-condizione} & L'utente ha inserito il comando \texttt{CONNECT}\\
				\noalign{\hrule height 0.5pt}
				\textbf{Post-condizione} & L'utente ha inserito con successo l'indirizzo del server \\
				\noalign{\hrule height 0.5pt}
				\textbf{Flusso eventi} & \begin{enumerate}
				\item L'utente scrive sul terminale l'indirizzo del server
				\end{enumerate} \\
				\noalign{\hrule height 0.5pt}
				\textbf{Scenari alternativi} & Nessuno \\
				\noalign{\hrule height 0.5pt}
				\textbf{Lista requisiti\newline dedotti} & \begin{itemize}
				\item ...
				\end{itemize} 
			\end{tabularx}
			\caption{Caso d'uso UC 1.2 - Inserimento indirizzo server}
		 \end{table} 
		 
		 
		 \subsection{UC 1.3: Inserimento username e password}
	 \textbf{Descrizione}
	 \\ \\
	 L'utente sta effettuando l'operazione di connessione al server e deve inserire il proprio username e la password di accesso al server.
	\begin{table}[H]
			\begin{tabularx}{\textwidth}{r  X}
				\textbf{Codice gerarchico} & UC1.3 \\
				\noalign{\hrule height 0.5pt}
				\textbf{Nome sintetico} & Inserimento username e password \\
				\noalign{\hrule height 0.5pt}
				\textbf{Attore principale} & Utente non autenticato\\
				\noalign{\hrule height 0.5pt}
				\textbf{Attori secondari} & Nessuno \\
				\noalign{\hrule height 0.5pt}
				\textbf{Pre-condizione} & L'utente ha inserito l'indirizzo del server\\
				\noalign{\hrule height 0.5pt}
				\textbf{Post-condizione} & L'utente ha inserito con successo username e password \\
				\noalign{\hrule height 0.5pt}
				\textbf{Flusso eventi} & \begin{enumerate}
				\item L'utente inserisce il proprio username
				\item L'utente inserisce la password 
				\end{enumerate} \\
				\noalign{\hrule height 0.5pt}
				\textbf{Scenari alternativi} & Nessuno \\
				\noalign{\hrule height 0.5pt}
				\textbf{Lista requisiti\newline dedotti} & \begin{itemize}
				\item ...
				\end{itemize} 
			\end{tabularx}
			\caption{Caso d'uso UC 1.3 - Inserimento username e password}
		 \end{table} 	
		 
		  
	 \subsection{UC 2: Visualizzazione aiuto}
	 \begin{figure}[H]
				\centering
				\includegraphics[scale=0.3]{UC/"UC 2 Visualizzazione aiuto".png}
				\caption{Diagramma di UC2: Visualizzazione aiuto}
			\end{figure}
	\textbf{Descrizione} 
	\\ \\
	L'utente (autenticato o non) intende ricevere un aiuto su come utilizzare i comandi di \emph{Actorbase}. Ha a disposizione due possibilità: 
	\begin{itemize}
	\item Aiuto semplice/generale
	\item Aiuto specifico
	\end{itemize}
	La prima modalità stampa sul terminale la lista completa dei comandi esistenti con relativa nota esplicativa sul loro comportamento. \\
	La modalità di aiuto specifica richiede l'inserimento del nome del comando per cui si richiedono le informazioni, stampa sul terminale la descrizione di tale comando.
	\begin{table}[H]
			\begin{tabularx}{\textwidth}{r X}
				\textbf{Codice gerarchico} & UC2 \\
				\noalign{\hrule height 0.5pt}
				\textbf{Nome sintetico} & Visualizzazione aiuto \\
				\noalign{\hrule height 0.5pt}
				\textbf{Attore principale} & Utente non autenticato e Utente autenticato\\
				\noalign{\hrule height 0.5pt}
				\textbf{Attori secondari} & Nessuno \\
				\noalign{\hrule height 0.5pt}
				\textbf{Pre-condizione} & L'utente ha avviato l'interfaccia CLI e intende visualizzare le informazioni di aiuto\\
				\noalign{\hrule height 0.5pt}
				\textbf{Post-condizione} & L'utente ha visualizzato le informazioni di aiuto richieste\\
				\noalign{\hrule height 0.5pt}
				\textbf{Flusso eventi} & \begin{enumerate}
				\item L'utente inserisce il comando relativo all'aiuto richiesto
				\item L'utente riceve le informazioni di aiuto stampate sul terminale
				\end{enumerate} \\
				\noalign{\hrule height 0.5pt}
				\textbf{Scenari alternativi} & Nessuno \\
				\noalign{\hrule height 0.5pt}
				\textbf{Lista requisiti\newline dedotti} & \begin{itemize}
				\item ...
				\end{itemize} 
			\end{tabularx}
			\caption{Caso d'uso UC 2 - Visualizzazione aiuto}
		 \end{table} 
		 
		 
		 \subsection{UC 2.1: Aiuto semplice}
	 \textbf{Descrizione}
	 \\ \\
	 L'utente sta effettuando l'operazione di visualizzazione dell'aiuto e intende inserire il comando di aiuto semplice.
	\begin{table}[H]
			\begin{tabularx}{\textwidth}{r  X}
				\textbf{Codice gerarchico} & UC2.1 \\
				\noalign{\hrule height 0.5pt}
				\textbf{Nome sintetico} & Aiuto semplice\\
				\noalign{\hrule height 0.5pt}
				\textbf{Attore principale} & Utente non autenticato e Utente autenticato\\
				\noalign{\hrule height 0.5pt}
				\textbf{Attori secondari} & Nessuno \\
				\noalign{\hrule height 0.5pt}
				\textbf{Pre-condizione} & L'utente intende ricevere un aiuto generale\\
				\noalign{\hrule height 0.5pt}
				\textbf{Post-condizione} & L'utente ha inserito correttamente il comando di aiuto semplice/generale\\
				\noalign{\hrule height 0.5pt}
				\textbf{Flusso eventi} & \begin{enumerate}
				\item L'utente inserisce il comando \texttt{HELP}
				\end{enumerate} \\
				\noalign{\hrule height 0.5pt}
				\textbf{Scenari alternativi} & Nessuno \\
				\noalign{\hrule height 0.5pt}
				\textbf{Lista requisiti\newline dedotti} & \begin{itemize}
				\item ...
				\end{itemize} 
			\end{tabularx}
			\caption{Caso d'uso UC 2.1 - Aiuto semplice}
		 \end{table} 	
	 
	 
	 \subsection{UC 2.2: Aiuto specifico}
	 \textbf{Descrizione}
	 \\ \\
	 L'utente sta effettuando l'operazione di visualizzazione dell'aiuto e intende inserire il comando di aiuto specifico.
	\begin{table}[H]
			\begin{tabularx}{\textwidth}{r  X}
				\textbf{Codice gerarchico} & UC2.2 \\
				\noalign{\hrule height 0.5pt}
				\textbf{Nome sintetico} & Aiuto specifico \\
				\noalign{\hrule height 0.5pt}
				\textbf{Attore principale} & Utente non autenticato e Utente autenticato\\
				\noalign{\hrule height 0.5pt}
				\textbf{Attori secondari} & Nessuno \\
				\noalign{\hrule height 0.5pt}
				\textbf{Pre-condizione} & L'utente intende ricevere un aiuto specifico\\
				\noalign{\hrule height 0.5pt}
				\textbf{Post-condizione} & L'utente ha inserito correttamente il comando di aiuto specifico\\
				\noalign{\hrule height 0.5pt}
				\textbf{Flusso eventi} & \begin{enumerate}
				\item L'utente inserisce il comando \texttt{HELP} seguito dal nome del comando per cui vuole ricevere l'aiuto
				\end{enumerate} \\
				\noalign{\hrule height 0.5pt}
				\textbf{Scenari alternativi} & Nessuno \\
				\noalign{\hrule height 0.5pt}
				\textbf{Lista requisiti\newline dedotti} & \begin{itemize}
				\item ...
				\end{itemize} 
			\end{tabularx}
			\caption{Caso d'uso UC 2.2 - Aiuto specifico}
		 \end{table} 	
	 
	 
	 \subsection{UC 3: Operazioni sui database}
	 \begin{figure}[H]
				\centering
				\includegraphics[scale=0.2]{UC/"UC 3 Operazioni sui database".png}
				\caption{Diagramma di UC3: Operazioni sui database}
			\end{figure}
	\textbf{Descrizione} 
	\\ \\
	L'utente ha effettuato correttamente la connessione, ha dunque accesso ai database presenti sul server. Su tali database può eseguire le seguenti operazioni:
	\begin{itemize}
	\item Visualizzare la lista dei database presenti
	\item Esportare uno o più database (backup)
	\item Importare un database
	\item Creare un database
	\item Eliminare un database
	\item Rinominare un database presente
	\item Selezionare un database per effettuare operazioni a livello di mappa su di esso
	\end{itemize}
	Queste operazioni possono portare a diverse situazioni di errore. In tal caso non viene eseguita alcuna operazione e l'utente riceve un messaggio di errore esplicativo.
	\begin{table}[H]
			\begin{tabularx}{\textwidth}{r X}
				\textbf{Codice gerarchico} & UC3 \\
				\noalign{\hrule height 0.5pt}
				\textbf{Nome sintetico} & Operazioni sui database\\
				\noalign{\hrule height 0.5pt}
				\textbf{Attore principale} & Utente autenticato\\
				\noalign{\hrule height 0.5pt}
				\textbf{Attori secondari} & Nessuno \\
				\noalign{\hrule height 0.5pt}
				\textbf{Pre-condizione} & L'utente si è connesso con successo ad un server\\
				\noalign{\hrule height 0.5pt}
				\textbf{Post-condizione} & L'operazione sui database selezionata è stata eseguita con successo\\
				\noalign{\hrule height 0.5pt}
				\textbf{Flusso eventi} & \begin{enumerate}
				\item L'utente inserisce il comando per l'operazione richiesta
				\end{enumerate} \\
				\noalign{\hrule height 0.5pt}
				\textbf{Scenari alternativi} & \begin{enumerate}
				\item L'utente ha richiesto un'operazione di esportazione e questa è fallita, l'operazione non viene eseguita e l'utente riceve un messaggio di errore (UC 3.8)
				\item L'utente ha richiesto un'operazione di importazione e questa è fallita, l'operazione non viene eseguita e l'utente riceve un messaggio di errore (UC 3.9)
				\item L'utente ha richiesto la creazione di un nuovo database sul server, la creazione è fallita e l'utente riceve un messaggio di errore (UC 3.10)
				\item L'utente ha tentato di accedere ad un database non presente sul server, l'accesso è fallito e l'utente riceve un messaggio di errore (UC 3.11)
\end{enumerate}				 \\
				\noalign{\hrule height 0.5pt}
				\textbf{Lista requisiti\newline dedotti} & \begin{itemize}
				\item ...
				\end{itemize} 
			\end{tabularx}
			\caption{Caso d'uso UC 3 - Operazioni sui database}
		 \end{table} 
	 
	 
	\subsection{UC 3.1: Visualizzazione lista dei database}
	\textbf{Descrizione} 
	\\ \\
	L'utente intende visualizzare la lista dei database presenti sul server a cui è connesso. Inserisce quindi il comando \texttt{SHOWDB} e il sistema stampa sul terminale la lista.
	\begin{table}[H]
			\begin{tabularx}{\textwidth}{r X}
				\textbf{Codice gerarchico} & UC3.1 \\
				\noalign{\hrule height 0.5pt}
				\textbf{Nome sintetico} & Visualizzazione lista dei database\\
				\noalign{\hrule height 0.5pt}
				\textbf{Attore principale} & Utente autenticato\\
				\noalign{\hrule height 0.5pt}
				\textbf{Attori secondari} & Nessuno \\
				\noalign{\hrule height 0.5pt}
				\textbf{Pre-condizione} & L'utente intende visualizzare la lista dei database presenti\\
				\noalign{\hrule height 0.5pt}
				\textbf{Post-condizione} & L'utente riceve la lista dei database presenti stampata a video\\
				\noalign{\hrule height 0.5pt}
				\textbf{Flusso eventi} & \begin{enumerate}
				\item L'utente inserisce il comando \texttt{SHOWDB}
				\end{enumerate} \\
				\noalign{\hrule height 0.5pt}
				\textbf{Scenari alternativi} & Nessuno \\
				\noalign{\hrule height 0.5pt}
				\textbf{Lista requisiti\newline dedotti} & \begin{itemize}
				\item ...
				\end{itemize} 
			\end{tabularx}
			\caption{Caso d'uso UC 3.1 - Visualizzazione lista dei database}
		 \end{table} 	 
	 
	 
	 
	 
	 
	\newpage \section{Requisiti}
		Di seguito vengono riportati tutti i requisiti individuati dal gruppo SWEeneyThreads. \\Tali requisiti
		derivano dai casi d’uso, dall'analisi del capitolato, dagli scambi di informazioni avvenuti con il
		\emph{proponente Cardin Riccardo}, oppure da necessità interne. \\
		La struttura di un requisito è definita nel documento \emph{Norme di progetto v1.1.1 sez 2.1.2}.
		I requisiti saranno elencati secondo un ordine. Ogni requisito seguirà la seguente codifica: \\
		\begin{center}
			R[Codice][Importanza][Tipo]
		\end{center}
		\textbf{Codice} \\ \\ Un codice univoco ed espresso in modo gerarchico;\\ \\
		\textbf{Importanza} \\ \\Può assumere i seguenti valori:
		\begin{itemize}
			\item \textbf{N:} necessario;
			\item \textbf{D:} desiderabile;
			\item \textbf{O:} opzionale.
		\end{itemize}
		\textbf{Tipo} \\ \\Può assumere i seguenti valori:
		\begin{itemize}
			\item \textbf{F:} funzionale;
			\item \textbf{Q:} di qualità;
			\item \textbf{P:} prestazionale;
			\item \textbf{V:} vincolo.
		\end{itemize}
	
	\subsection{Requisiti funzionali}
	I requisiti funzionali riguardano le funzioni vere e proprie del prodotto. Il soddisfacimento di uno di essi
	equivale all'implementazione di una funzionalità.
	\LTXtable{\textwidth}{tabelle_requisiti/funzionali.tex}
	\subsection{Requisiti di vincolo}
	I requisiti di vincolo rappresentano i vincoli che devono essere soddisfatti dal prodotto e che esulano dalle
	sue caratteristiche funzionali.
	\LTXtable{\textwidth}{tabelle_requisiti/vincolo.tex}
	\subsection{Requisiti di qualità}
	I requisiti di qualità specificano le operazioni da compiere per far raggiungere al prodotto il livello 
	qualitativo richiesto.
	\LTXtable{\textwidth}{tabelle_requisiti/qualita.tex}
	\subsection{Requisiti prestazionali}
	I requisiti prestazionali consentono, al loro soddisfacimento, di far raggiungere al prodotto un determinato
	livello di prestazioni, non si traducono in nuove funzionaltà per l'utente ma intendono migliorare 
	la stabilità e la velocità del programma.
	\LTXtable{\textwidth}{tabelle_requisiti/prestazionali.tex}
	
	\subsection{Tracciamento fonti-requisiti}
	Di seguito si riporta in forma tabellare l'elenco di requisiti ricavati dalle diverse fonti.
	\LTXtable{\textwidth}{tabelle_requisiti/fontirequisiti.tex}
	
	\cleardoublepage
	\addcontentsline{toc}{section}{\listfigurename}
	\listoffigures
	
	\cleardoublepage
	\addcontentsline{toc}{section}{\listtablename}
	\listoftables
		
\end{document}
	
